\par\medskip\noindent
\textbf{\SGAInlineRefTarget{sga3:xix:paragraph:3-1}3.1.}---
Let \(S\) be a scheme, and let \(T\) be an \(S\)-torus acting linearly
on a finite locally free \(\mathcal O_S\)-module \(F\) (cf. I, §\SGARefLink{sga3:xix:lemma:4-7}{4.7}).
For every character \(\alpha\) of \(T\), that is,
\[
  \alpha\in
  \operatorname{Hom}_{S\text{-gr.}}(T,\mathbb G_{m,S}),
\]
define a subfunctor of \(\mathbb W(F)\) by
\[
  \mathbb W(F)^\alpha(S')
  =
  \left\{
    x\in\mathbb W(F)(S')
    \ \middle|\
    t\cdot x=\alpha(t)x
    \text{ for every }t\in T(S''),\ S''\to S'
  \right\}.
\]
\SGARefTarget{sga3:xix:paragraph:3:1}

\par\medskip\noindent
\textbf{\textit{Lemma.}}---%
\footnote{Editor's note: this statement has been made into a lemma in
order to bring it out.}
\textit{Then
\(\mathbb W(F)^\alpha=\mathbb W(F^\alpha)\), where \(F^\alpha\) is a
submodule of \(F\), locally a direct summand of \(F\), and hence also
locally free.}
\SGARefTarget{sga3:xix:lemma:3:1}

Indeed, the assertion is local for the fpqc topology, and we may
suppose that \(T=D_S(M)\), where \(M\) is a finite-type (free) abelian
group.  Then \(\alpha\) is identified with a locally constant function
from \(S\) to \(M\) (\SGARefLink{sga3:viii:corollary:1:3}{Exposé VIII, 1.3}), and, after restricting \(S\) if
necessary, we may suppose that this function is constant.  We are then
reduced to \SGARefLink{sga3:i:proposition:4-7-3}{Exposé I, 4.7.3}.

\par\medskip\noindent
\SGAInlineRefTarget{sga3:xix:definition:3-2}\textbf{\textit{\SGARefLink{sga3:xix:definition:3:2}{Definition 3.2}.}}---
\textit{Let \(S\) be a scheme, let \(G\) be an \(S\)-group scheme
smooth over \(S\) and with connected fibers,\footnote{Editor's note:
the original also assumed that \(G\) was of finite presentation over
\(S\), an assumption that does not appear to be used below.  In any
case, since \(G\) is smooth over \(S\) and has connected fibers, it is
quasi-compact and separated over \(S\) (VI\textsubscript{B}, \SGARefLink{cor:sga3-VIB-5.5}{5.5}), and
hence of finite presentation over \(S\).} and let \(T\) be a torus of
\(G\).  Put \(\mathfrak g=\operatorname{Lie}(G/S)\), and let \(T\)
act on \(\mathfrak g\) through the adjoint representation of \(G\).}

\textit{A character \(\alpha\) of \(T\) is called a \emph{root} of
\(G\) with respect to \(T\) if the following equivalent conditions
are satisfied:}
\SGARefTarget{sga3:xix:definition:3:2}

\begin{enumerate}
\item[\textit{(i)}]
\textit{For every \(s\in S\), \(\alpha_s\) is a root of \(G_s\) with
respect to \(T_s\) (\SGARefLink{sga3:xix:paragraph:1-10}{1.10}).}

\item[\textit{(ii)}]
\textit{\(\alpha\) is nontrivial on every fiber, and
\(\mathfrak g^\alpha(s)\ne0\) for every \(s\in S\).}
\end{enumerate}

\par\medskip\noindent
\textbf{\textit{\SGARefLink{sga3:xix:lemma:3:3}{Lemma 3.3}.}}---
\textit{Let \(S\) be a scheme, let \(T\) be an \(S\)-torus, and let
\(\alpha\) be a character of \(T\).  The following conditions are
equivalent:}
\SGARefTarget{sga3:xix:lemma:3:3}

\begin{enumerate}
\item[\textit{\SGAInlineRefTarget{sga3:xix:lemma:3:3:item:i}(i)}]
\textit{\(\alpha\) is nontrivial on every fiber; that is, for every
\(s\in S\), \(\alpha_s\) is distinct from the unit character of
\(T_s\).}

\item[\textit{\SGAInlineRefTarget{sga3:xix:lemma:3:3:item:ii}(ii)}]
\textit{For every \(S'\to S\) with \(S'\ne\varnothing\),
\(\alpha_{S'}\) is distinct from the unit character of \(T_{S'}\).}

\item[\textit{\SGAInlineRefTarget{sga3:xix:lemma:3:3:item:iii}(iii)}]
\textit{The morphism \(\alpha\) is faithfully flat.}
\end{enumerate}

\footnote{Editor's note: the original has been simplified; it invoked
\SGARefLink{sga3:ix:corollary:5:2}{Exposé IX, 5.2}.  Note, however, that loc.\ cit., \SGARefLink{sga3:xix:lemma:5-3}{5.3} shows that, if
\(S\) is connected, the conditions of the lemma hold as soon as
\textup{\SGARefLink{sga3:xix:lemma:3:3:item:i}{(i)}} holds at one point \(s\) of \(S\).}
It is clear that \SGARefLink{sga3:xix:lemma:3:3:item:ii}{(ii)} implies \SGARefLink{sga3:xix:lemma:3:3:item:i}{(i)}, and it is easy to see that \SGARefLink{sga3:xix:lemma:3:3:item:iii}{(iii)}
implies \SGARefLink{sga3:xix:lemma:3:3:item:i}{(i)}.  Condition \SGARefLink{sga3:xix:lemma:3:3:item:i}{(i)} implies \SGARefLink{sga3:xix:lemma:3:3:item:ii}{(ii)}: if \(s'\in S'\) lies over
\(s\), and \(\alpha_{s'}\) is the unit character, then so is
\(\alpha_s\).  Finally, to prove that \SGARefLink{sga3:xix:lemma:3:3:item:i}{(i)} implies \SGARefLink{sga3:xix:lemma:3:3:item:iii}{(iii)}, we reduce to
the case in which \(T\) is diagonalizable and conclude by \SGARefLink{sga3:viii:corollary:3:2:item:a:01}{Exposé VIII,
3.2(a)}.

\par\medskip\noindent
\textbf{3.4.}---%
\footnote{Editor's note: what follows has been modified so as to recall
the hypotheses of \SGARefLink{sga3:xix:definition:3-2}{3.2} and to add that \(T\) is a maximal torus.}
Let \(G\) be a reductive \(S\)-group scheme, let \(T\) be a maximal
torus of \(G\), and let \(\alpha\) be a root of \(G\) with respect to
\(T\).  Then, by \SGARefLink{sga3:xix:notations:2:5:0}{2.5.0} and \SGARefLink{sga3:xix:proposition:1-12}{1.12}, \(\mathfrak g^\alpha\) is a locally
free \(\mathcal O_S\)-module of rank one.  Moreover, by \SGARefLink{sga3:xix:paragraph:1-10}{1.10},
\(-\alpha\) is also a root of \(G\) with respect to \(T\).  In
particular, if \(G\) has semisimple rank one, \SGARefLink{sga3:xix:lemma:1-11}{1.11} and \SGARefLink{sga3:xix:proposition:1-12}{1.12} give the
following statement.
\SGARefTarget{sga3:xix:paragraph:3:4}
